% SOURCE_AUTHORITY: sga5_fr_workpass.tex, lines 11669--11727 (LF-normalized unit).
% SOURCE_SHA256: 791F4EFFC5E02832D5D77ED03518C8156D6F07E4C8238B03545DB93D883FBB28
\subsection*{2. Caso de categorias trianguladas}

Seja $\cC$ uma categoria triangulada e $G$ um grupo abeliano. Diz-se que uma função $f$ de $\cC$ em $G$ é aditiva se, para todo triângulo distinguido
\[
\begin{tikzcd}[column sep=small,row sep=small]
& Z \arrow[dl,"w"']\\
X \arrow[rr,"u"'] && Y \arrow[ul,"v"']
\end{tikzcd}
\]
tem-se
\[
f(Y)=f(X)+f(Z).
\]

Dada uma função aditiva $f$, as seguintes relações são consequências imediatas da definição:
\[
\text{a) } f(0)=0,\qquad
\text{b) } f(X[n])=(-1)^n f(X),
\]
\[
\text{c) } f(X)=f(Y)\text{ se }X\text{ e }Y\text{ são isomorfos,}
\qquad
\text{d) } f(X\oplus Y)=f(X)+f(Y).
\]
Os triângulos seguintes são distinguidos:
\[
\begin{tikzcd}[column sep=small,row sep=small]
&0 \arrow[dl]\\ X \arrow[rr,"1_X"']&&X\arrow[ul]
\end{tikzcd}
\qquad
\begin{tikzcd}[column sep=small,row sep=small]
&X[1]\arrow[dl]\\ X\arrow[rr]&&0\arrow[ul]
\end{tikzcd}
\]
\[
\begin{tikzcd}[column sep=small,row sep=small]
&0\arrow[dl]\\ X\arrow[rr,"u"']&&Y\arrow[ul]
\end{tikzcd}
\qquad
\begin{tikzcd}[column sep=small,row sep=small]
&Y[1]\arrow[dl]\\ X\oplus Y\arrow[rr]&&X\arrow[ul]
\end{tikzcd}
\]
($u$ isomorfismo). Para formar o último triângulo, utilizamos o fato de que a soma de dois triângulos distinguidos é distinguida.

Como no caso das categorias abelianas, obtém-se evidentemente um funtor
\[
A:\Ab\longrightarrow\Ab,
\]
que associa a cada grupo abeliano $G$ o grupo abeliano $A(G)$ das funções aditivas sobre $\cC$ com valores em $G$. Este funtor é representável e é fácil construir um par de representação $(K(\cC),\operatorname{cl}_{\cC})$. Com efeito, considera-se primeiro o grupo abeliano livre $\mathbb Z^{(\Ob\cC)}$ gerado pelos objetos de $\cC$ e, depois, o subgrupo $R$ gerado pelos elementos da forma $X-Y-Z$, onde
\[
\begin{tikzcd}[column sep=small,row sep=small]
&Z\arrow[dl]\\ Y\arrow[rr]&&X\arrow[ul]
\end{tikzcd}
\]
é um triângulo distinguido. O grupo quociente $\mathbb Z^{(\Ob\cC)}/R$ será $K(\cC)$, e $\operatorname{cl}_{\cC}$ é obtido compondo as duas aplicações canônicas:
\[
\Ob\cC\longrightarrow \mathbb Z^{(\Ob\cC)}\longrightarrow K(\cC).
\]
